\documentclass[aps,prl,onecolumn,superscriptaddress,nofootinbib]{revtex4-2}

\usepackage{amsmath,amssymb,amsthm,mathtools,bm}
\usepackage{booktabs}
\usepackage{graphicx}
\usepackage{xcolor}
\usepackage[colorlinks=true,linkcolor=blue,citecolor=blue,urlcolor=blue]{hyperref}
\usepackage{microtype}

\renewcommand{\thesection}{S\Roman{section}}
\renewcommand{\theequation}{S\arabic{equation}}
\renewcommand{\thefigure}{S\arabic{figure}}
\renewcommand{\thetable}{S\arabic{table}}

\newcommand{\Tr}{\operatorname{Tr}}
\newcommand{\avg}[1]{\overline{#1}}
\newcommand{\Gtraj}{G_{\xi}}
\newcommand{\Gbar}{\overline{G}}
\newcommand{\Gps}{G_{\rm ps}}
\newcommand{\GOW}{G_{\rm OW}}
\newcommand{\calL}{\mathcal{L}}
\newcommand{\calC}{\mathcal{C}}
\newcommand{\calD}{\mathcal{D}}
\newcommand{\calE}{\mathcal{E}}
\newcommand{\calN}{\mathcal{N}}
\newcommand{\dd}{\mathrm{d}}
\newcommand{\nsh}{n_{\rm shell}}
\newcommand{\atopo}{\alpha_{\rm top}}
\newcommand{\atriv}{\alpha_{\rm triv}}
\newcommand{\Dlam}{\Delta_{\lambda}}
\newcommand{\Dp}{\Delta_{\rm p}}
\newcommand{\ket}[1]{|#1\rangle}
\newcommand{\bra}[1]{\langle#1|}
\newcommand{\braket}[2]{\langle#1|#2\rangle}

\newcommand{\figplaceholder}[2]{%
  \fbox{\begin{minipage}[c][#1][c]{0.96\linewidth}
    \centering\small\textsf{#2}
  \end{minipage}}}

\begin{document}

\title{Supplemental Material for ``Steering fermions to chiral criticality''}

\author{Asadullah Bhuiyan}
\affiliation{Department of Physics, Cornell University, Ithaca, New York 14853, USA}
\author{Haining Pan}
\affiliation{Department of Physics and Astronomy, Center for Materials Theory, Rutgers University, Piscataway, New Jersey 08854, USA}
\author{Chao-Ming Jian}
\affiliation{Department of Physics, Cornell University, Ithaca, New York 14853, USA}

\date{\today}

\maketitle
\setcounter{tocdepth}{2}
\tableofcontents

%=====================================================================
\section{Model, overcomplete-Wannier construction, and circuit algorithm}
\label{sm:model}

\subsection{Parent Hamiltonian and domain-wall geometry}

We consider spinless complex fermions $\hat c_{\bm r,\mu}$, $\mu=1,2$, on an $N_x\times N_y$ square lattice with periodic boundary conditions in both directions, together with an identical ancillary layer $\hat d_{\bm r,\mu}$ playing the role of a particle bath.  The target-state structure is generated by the two-band Chern-insulator parent Hamiltonian
\begin{equation}
  h(\bm k)=\sin k_x\,\sigma^x+\sin k_y\,\sigma^y+\big(\alpha-\cos k_x-\cos k_y\big)\sigma^z,
  \label{eq:sm-parent}
\end{equation}
whose lower band carries Chern number $-1$ for $\alpha\in(0,2)$, $+1$ for $\alpha\in(-2,0)$, and $0$ for $|\alpha|>2$ \cite{BhuiyanPanJian2026}.  The domain-wall geometry is defined by a position-dependent mass: $\alpha(x)=\atopo$ in the topological slab and $\alpha(x)=\atriv$ in the trivial exterior.  Unless stated otherwise, production runs use $\atopo=1$, $\atriv=30$, walls at $x=6$ and $x=14$ for $N_x=20$, so that the slab $x\in[6,14]$ is topological and the exterior is strongly trivial.  Because the mass conventions differ across internal data sets, we consistently use the physical labels ``topological mass'' $\atopo$ and ``trivial mass'' $\atriv$.

\subsection{Overcomplete-Wannier modes and stabilizer conditions}

Chern bands admit no complete basis of exponentially localized Wannier functions, but they do admit an \emph{overcomplete} localized basis \cite{BhuiyanPanJian2026,Rashba1997,Qi2011,JianQi2013}.  For each unit cell $\bm r$, trial orbital $\nu=A,B$ with spinors $\tau_A=(1,1)^{\sf T}/\sqrt2$, $\tau_B=(1,-1)^{\sf T}/\sqrt2$, and band index $\pm$, the localized OW fermion operator is
\begin{equation}
  \hat\chi_{\bm r,\nu,\pm}\propto\int\!\frac{\dd^2\bm k}{(2\pi)^2}\,
  e^{\mathrm{i}\bm k\cdot\bm r}\,\tau^\dagger_{\nu}P_{\pm}(\bm k)\,\hat c(\bm k),
  \qquad
  P_\pm(\bm k)=\tfrac12\Big(\mathbf{1}\pm\tfrac{\vec n(\bm k)\cdot\vec\sigma}{|\vec n(\bm k)|}\Big),
  \label{eq:sm-ow}
\end{equation}
with $\vec n(\bm k)$ the Bloch vector of Eq.~\eqref{eq:sm-parent}.  In the domain-wall geometry the projectors are built from the real-space Hamiltonian with the mass profile $\alpha(x)$.  The OW wavefunctions are exponentially localized whenever the local gap is finite ($\alpha\neq0,\pm2$).  The target state $\ket{\rm CI}$ is the unique common eigenstate of the stabilizing conditions
\begin{equation}
  \hat\calN_{\bm r,\nu,-}\ket{\rm CI}=\ket{\rm CI},
  \qquad
  \big(1-\hat\calN_{\bm r,\nu,+}\big)\ket{\rm CI}=\ket{\rm CI},
  \qquad
  \hat\calN_{\bm r,\nu,\pm}\equiv\hat\chi^\dagger_{\bm r,\nu,\pm}\hat\chi_{\bm r,\nu,\pm},
  \label{eq:sm-stab}
\end{equation}
for all $\bm r,\nu$; the overcompleteness of the OW basis guarantees uniqueness \cite{BhuiyanPanJian2026}.

\subsection{Adaptive cycle and finite-range truncation}

One adaptive cycle iterates over all unit cells and OW labels: (i) measure $\hat\calN_{\bm r,\nu,\pm}$ (Born-sampled); (ii) if the outcome violates Eq.~\eqref{eq:sm-stab}, apply the on-site error-correcting fermionic SWAP ${\rm fSWAP}(\hat\chi_{\bm r,\nu,\pm},\hat d_{\bm r,\nu})$ that exchanges the offending particle (hole) with the ancillary mode; (iii) after the sweep, randomly redistribute charge in the ancillary layer by charge-conserving Gaussian unitaries and measure all ancillary occupations.  The feedback is deterministic in its \emph{target} but succeeds with the finite probability that the ancillary mode can absorb or supply a particle; this is the ``perfect correction'' protocol.  All operations are strictly local and require no classical communication.

Finite-range truncation keeps each OW mode on a $(2\nsh{+}1)\times(2\nsh{+}1)$ neighborhood of its center (the ``$\nsh$-shell'' square mask), followed by renormalization.  Truncated number operators on disjoint supports commute, so measurement-feedforward operations can be grouped into a lattice-size-independent number of parallel layers: the circuit depth per cycle is $\mathcal{O}(1)$ at fixed $\nsh$ \cite{BhuiyanPanJian2026}.  The additional \texttt{dw\_truncation} mask restricts each OW mode to the sector (topological or trivial) containing its center; its interpretation is developed in Sec.~\ref{sm:owtrunc}.

\subsection{Numerical engines and simulation parameters}

All trajectory simulations call the canonical adaptive Markov-circuit drivers of the internal code base---\texttt{classA\_U1FGTN.run\_markov\_circuit} (CPU) and \texttt{classA\_U1FGTN\_gpu.run\_markov\_circuit} (GPU)---which propagate the top-layer covariance $\Gtraj=2C_\xi-\mathbf{1}$ exactly through Born-sampled measurements and feedback, using a rank-one resolvent update per elementary event.  Trajectory-averaged channel evolution uses the separate deterministic driver \texttt{run\_markov\_channel} (Sec.~\ref{sm:avg}).  Table~\ref{tab:params} summarizes the main production campaigns used in the Letter.

\begin{table}[h]
\caption{\label{tab:params}Principal simulation campaigns.  PC = perfect correction; PS = full postselection (deterministic forced branch; a single trajectory by construction).  All campaigns use $\nsh=1$, trial orbitals $A,B$ (or $X$ in legacy labeling), $\atopo=1$, $\atriv=30$, and the hard-wall (\texttt{dw\_truncation}) construction unless stated.}
\begin{ruledtabular}
\begin{tabular}{llllll}
Campaign & Geometry & Initial state & Cycles & Samples & Used in \\
\colrule
Max-mixed purification (GPU) & $20\times\{30,40,50\}$ & maximally mixed & $2N_y$ & 100 (PC), 1 (PS) & Fig.~2(a,b,d) \\
Charge-sharpening $\alpha$ sweep (CPU) & $16\times\{16,24,32\}$ & maximally mixed & $2N_y$ & 10 (PC), 1 (PS) & Fig.~2(c) \\
Large-$N_y$ entanglement scaling (GPU) & $20\times\{40,\dots,120\}$ & domain-wall default & $N_y/2$ & 10 & Fig.~3(a,b) \\
Fixed-cycle size sweep (GPU) & $20\times\{30,40,50\}$ & domain-wall default & 50 & 100 & Fig.~3(b) \\
Streaming characterization (GPU) & $20\times\{30,40,50\}$ & domain-wall default & $2N_y$ & 100 (PC), 1 (PS) & Fig.~3(c), Sec.~\ref{sm:corr} \\
Modular charge spreading (CPU) & $20\times40$ & prepared/exact states & modular time & --- & Fig.~3(d), Sec.~\ref{sm:modular} \\
Lyapunov spectra (GPU) & $20\times\{30,40,50\}$ & domain-wall default & $N_y$ & 25 (PC), 1 (PS) & Fig.~4(a,b,d) \\
Choi transfer, CPU & $20\times\{20,30,40\}$ & domain-wall default & $2N_y$ & 10 & Fig.~4(c) \\
Small-system purification (GPU) & $16\times32$, $30\times40$ & maximally mixed & 100 & 10 & Sec.~\ref{sm:owtrunc} \\
Partial postselection (GPU) & $20\times40$ & domain-wall default & 40 & 10 ($p<1$), 1 ($p=1$) & Sec.~\ref{sm:taxonomy} \\
\end{tabular}
\end{ruledtabular}
\end{table}

%=====================================================================
\section{Protocol taxonomy and ensemble ordering}
\label{sm:taxonomy}

\subsection{Four dynamical objects}

The elementary measurement-feedback event is a state-dependent random map: given the covariance $G$, the Born probability of each branch is computed from $G$, a branch is sampled, and the corresponding (deterministic) projection-plus-feedback map is applied.  The trajectory covariance $\Gtraj(t)$ is therefore a Markov process on covariance matrices, labeled by the record $\xi$.  Four related but inequivalent dynamical descriptions appear in this work:

\begin{enumerate}
\item \emph{Perfect correction (the physical protocol).}  Born outcomes are sampled; feedback toward the stabilizer target is applied deterministically upon a mismatch.  Observables are ensemble averages $\avg{O[\Gtraj]}=\sum_\xi p_\xi\,O[G_\xi]$ over corrected Born trajectories.  ``Perfect'' refers to the classical correction step, \emph{not} to selecting measurement outcomes.

\item \emph{Full postselection (control).}  Every measurement outcome is forced to its target value; no Born sampling occurs at the measured site, and the trajectory is deterministic.  The numerical drivers accordingly store a single sample for postselected runs regardless of the requested ensemble size; postselected entries must never be ensemble-weighted against perfect-correction entries.  Physically, the forced branch has weight exponentially small in spacetime volume \cite{LiZou2023,GarrattAltman2024}.

\item \emph{Partial postselection (interpolation control).}  An external Bernoulli mask with probability $p$ forces a site's outcome; the complement is Born-sampled with feedback.  In the production sweep at $20\times40$, $40$ cycles, $p\in\{0,0.1,0.25,0.5,0.75,0.9,0.95,0.99,1\}$ with $10$ samples for $p<1$, the late-time per-cycle covariance motion (Frobenius norm of successive differences) decreases smoothly from $\simeq11.0$ at $p=0$ to $\simeq2.2$ at $p=0.99$, collapsing to $\simeq0.06$ at the deterministic endpoint $p=1$; the endpoint also shows a degraded real-space Chern marker ($\simeq0.77$ vs.\ $\simeq1$ for $p<1$ samples).  These data demonstrate a smooth operational interpolation at fixed finite size; they are protocol diagnostics, not a thermodynamic phase diagram.

\item \emph{Mean channel and Lindblad limit (controls).}  The Born-weighted average $\Gbar=\avg{\Gtraj}$ evolves under a deterministic channel; its continuous-time embedding is a Lindbladian.  These objects are analyzed in Sec.~\ref{sm:avg}.
\end{enumerate}

\begin{figure}[h]
  \figplaceholder{1.5in}{%
  \textbf{FIG.~S1 PLACEHOLDER (partial postselection, $1\times3$ panels).}\\[2pt]
  (a) Late-time per-cycle covariance motion (Frobenius successive delta) vs.\ forced-measurement probability $p$, $20\times40$, 40 cycles: $\simeq11.0$ ($p=0$) $\to$ $\simeq2.2$ ($p=0.99$) $\to$ $\simeq0.06$ ($p=1$); (b) final-sample real-space Chern marker vs.\ $p$ (endpoint drop to $\simeq0.77$); (c) measurement-frustration payload vs.\ $p$ ($p<1$ only).\\
  {[}data: \texttt{colab\_partial\_post-select} run index and scalar metrics{]}}
  \caption{Partial-postselection interpolation control between the Born-sampled perfect-correction ensemble ($p=0$) and the deterministic postselected branch ($p=1$).}
  \label{fig:sm-partialps}
\end{figure}

\subsection{Ensemble ordering for nonlinear observables}

For a linear one-body observable $A$, averaging commutes with evaluation, $\avg{\Tr(AC_\xi)}=\Tr(A\bar C)$.  For nonlinear observables it does not:
\begin{equation}
  \avg{O[\Gtraj]}\neq O[\Gbar]\neq O[\Gps].
  \label{eq:sm-taxonomy}
\end{equation}
The cleanest exact statement is the variance decomposition of the squared correlator.  Writing $C_\xi(r,r')=\bar C(r,r')+\delta C_\xi(r,r')$ with $\avg{\delta C_\xi}=0$,
\begin{equation}
  \avg{|C_\xi(r,r')|^2}
  =|\bar C(r,r')|^2+\avg{|\delta C_\xi(r,r')|^2}.
  \label{eq:sm-decomp}
\end{equation}
The second term is the trajectory-to-trajectory fluctuation contribution, which any mean-channel or Lindblad treatment deletes by construction.  Entanglement entropies, entanglement contours, and regularized Chern markers are likewise nonlinear and do not commute with trajectory averaging.  Empirically the averaging order is \emph{not} responsible for the small residual entropy-slope excess of the Letter: the mean of per-sample slopes and the slope of the mean entropy curve agree to numerical roundoff (Sec.~\ref{sm:corr}).

%=====================================================================
\section{Trajectory-averaged evolution: mean channel, Lindbladian, and the mixed-state interface sector}
\label{sm:avg}

\subsection{Discrete-time mean channel}

Averaging the elementary measurement-feedback event over Born outcomes defines a deterministic, nonlinear map on the covariance,
\begin{equation}
  \Gbar_{n+1}=\Phi(\Gbar_n),
  \label{eq:sm-channel}
\end{equation}
implemented by the deterministic driver \texttt{classA\_U1FGTN.run\_markov\_channel}.  This computes $O[\Gbar]$, not $\avg{O[\Gtraj]}$; per Eq.~\eqref{eq:sm-decomp}, the two differ for every nonlinear $O$.  Every trajectory-averaged adaptive circuit corresponds to a quantum channel and hence---upon a continuous-time embedding---to a Lindbladian \cite{BhuiyanPanJian2026,Guo2025}.

\subsection{Lindblad embedding and measurement-induced interactions}

Assigning each elementary measurement-feedforward channel an infinitesimal rate embeds the cycle channel into a continuous-time Lindbladian.  Within the Markovian-bath approximation for the ancillary layer (fixed local ancillary density $\bar n_{\rm a}$), the generator obtained in Ref.~\cite{BhuiyanPanJian2026} is
\begin{equation}
  \calL^{\rm cycle}=\calL^{\rm gain}+\calL^{\rm loss}+\calL^{\rm dephas.},
  \label{eq:sm-lindblad}
\end{equation}
\begin{equation}
  \calL^{\rm gain}=\bar n_{\rm a}\!\sum_{\bm r,\nu}\calD[\hat\chi^\dagger_{\bm r,\nu,-}],\qquad
  \calL^{\rm loss}=(1-\bar n_{\rm a})\!\sum_{\bm r,\nu}\calD[\hat\chi_{\bm r,\nu,+}],\qquad
  \calL^{\rm dephas.}=(2-\bar n_{\rm a})\!\sum_{\bm r,\nu}\calD[\hat\calN_{\bm r,\nu,-}]
  +(1+\bar n_{\rm a})\!\sum_{\bm r,\nu}\calD[\hat\calN_{\bm r,\nu,+}],
  \label{eq:sm-lindblad-terms}
\end{equation}
with $\calD[\hat L](\hat\rho)=\hat L\hat\rho\hat L^\dagger-\tfrac12\{\hat L^\dagger\hat L,\hat\rho\}$.  The gain and loss jump operators are linear in fermions and preserve Gaussianity; the dephasing jump operators $\hat\calN_{\bm r,\nu,\pm}$ are \emph{quartic}.  Consequently, although every individual trajectory is a pure Gaussian state, the trajectory-averaged density matrix is generically non-Gaussian, and the averaged dynamics is that of an \emph{interacting} open system \cite{Dolgirev2020,BarthelZhang2022}.  This places the averaged sector outside the quadratic-Lindbladian symmetry classification \cite{LieuMcGinleyCooper2020,AltlandFleischhauerDiehl2021}.  Two structural facts make it nonetheless tractable:
(i) because the dephasing jump operators are Hermitian fermion bilinears, $2k$-point functions close on $\leq2k$-point functions, so the two-point function $\Gbar$ obeys a closed (nonlinear) equation of motion \cite{BhuiyanPanJian2026,BarthelZhang2022}; and
(ii) the closed two-point dynamics admits a vectorized affine representation $\frac{\dd}{\dd t}\mathrm{vec}(G)=s+L\,\mathrm{vec}(G)$ in the weak-update regime, implemented by the internal Lindblad helper (\texttt{CI\_Lindblad\_DW.py}).
For the uniform topological geometry, the dissipation gaps $\gamma_\pm(\bm k)=\sum_\nu|f_{\nu,\pm}(\bm k)|^2$ are strictly positive and $\mathcal{O}(1)$---the overcompleteness of the OW basis patches the topological zeros of the individual form factors---yielding the $T_{\rm conv}\sim\mathcal{O}(1)$ convergence of the bulk protocol \cite{BhuiyanPanJian2026}.

\subsection{Goldstein's no-go theorem and how the trajectory ensemble evades it}

At the averaged level, the steering protocol is a finite-range dissipative dynamics, and Goldstein's no-go theorem \cite{Goldstein2019} applies: no finite-range Lindbladian in $d\geq2$ can relax at a finite, size-independent rate to a \emph{unique pure} topological steady state; finite-range dissipators can only approach such a state through mixed states arbitrarily close to it.  Consistently, the averaged steady state of our domain-wall dynamics is mixed.  The trajectory-resolved protocol circumvents the obstruction because purity is restored at the level of individual Born trajectories: the stochastic measurement record performs the localization that no finite-range deterministic dissipator can, and the chiral critical correlations of the Letter live entirely in the trajectory ensemble.  This division of labor---deterministic skeleton in $\Gbar$, criticality in the fluctuations---is the operational content of the dichotomy stated in the main text.

\subsection{The mixed-state interface sector}
\label{sm:avg-interface}

In the domain-wall geometry the averaged steady state retains a coherent interface structure, summarized by four diagnostics (Fig.~\ref{fig:sm-avg}):
(i) the $k_y$-resolved occupation spectrum of $\Gbar$ exhibits branches interpolating between the near-filled and near-empty bulk sectors---\emph{edge bands} of the mixed state---whose eigenvectors are localized on the walls;
(ii) the covariance eigenmodes nearest to half filling are wall-localized;
(iii) the local Chern marker organizes into a strip pattern following the programmed mass profile;
(iv) current-like response is concentrated on the walls.
We deliberately phrase these as a ``mixed-state interface sector'' rather than a quantized mixed-state invariant; establishing a quantized invariant for this interacting averaged dynamics would require a separate definition and robustness analysis (cf.\ \cite{BudichDiehl2015,Bardyn2018,WawerFleischhauer2021}).  These observations are the internal-interface, interacting-channel counterparts of edge structure known for engineered quadratic dissipators at lattice boundaries \cite{Diehl2011,Bardyn2013,BudichZollerDiehl2015,YangMoligniniBergholtz2023,HegdeEhmckeMeng2023,ShavitGoldstein2020}.  Two cautions: discrete mean-channel data at update probability one are not by themselves a controlled continuous-time extrapolation, and the panels of Fig.~\ref{fig:sm-avg} are to be regenerated from the canonical deterministic drivers for archival provenance.

\begin{figure}[h]
  \figplaceholder{2.2in}{%
  \textbf{FIG.~S2 PLACEHOLDER (trajectory-averaged sector, up to $3\times2$ panels).}\\[2pt]
  (a) Mean-channel covariance flow: distance of $\Gbar_n$ to its fixed point vs.\ cycle, uniform vs.\ domain-wall geometry.  (b) $k_y$-resolved occupation spectrum of the steady $\Gbar$ with wall-localized edge branches; (c) spatial density of the mid-gap covariance eigenmodes; (d) local Chern-marker map of $\Gbar$; (e) wall-centered current-like response; (f) observable comparison: a linear observable where $\avg{O[\Gtraj]}=O[\Gbar]$ exactly, and entropy/$|C|^2$ where they split [Eq.~(\ref{eq:sm-decomp})].\\
  {[}to generate: \texttt{classA\_U1FGTN.run\_markov\_channel}, \texttt{CI\_Lindblad\_DW.py}, plus trajectory data from the streaming campaign{]}}
  \caption{The trajectory-averaged (mean-channel/Lindblad) sector in the domain-wall geometry.}
  \label{fig:sm-avg}
\end{figure}

%=====================================================================
\section{Exact domain-wall benchmark}
\label{sm:exact}

The exact reference state is the ground state of the real-space domain-wall Hamiltonian built from Eq.~\eqref{eq:sm-parent} at $N_x=20$, $N_y=40$, $\atopo=1$, $\atriv=30$, walls at $x=6,14$, periodic boundary conditions: diagonalizing the one-body Hamiltonian, occupying all negative-energy modes, and forming $G=2P_--\mathbf{1}$.  The single-particle spectrum contains $800$ negative and $800$ positive modes with finite-size minimum gap $|E|_{\min}=7.3\times10^{-15}$, reflecting the chiral crossings at the two walls.

\emph{Correlators.}  For a chiral free fermion on a ring of circumference $N_y$, $|\langle\psi^\dagger(y)\psi(0)\rangle|^2\sim[\tfrac{N_y}{\pi}\sin(\tfrac{\pi y}{N_y})]^{-2}$.  The exact state reproduces this with high precision and sharp spatial localization: log-chord fits of the squared correlator give slope $-2.0246$ on the wall columns $x=6,14$ ($R^2=0.99984$) and $-2.0689$ on the trivial-side neighbors $x=5,15$ ($R^2=0.99878$), while interior columns $x=7,13$ fall steeply with effective slope $\simeq-13.1$.  The $x$-averaged squared correlator remains edge-dominated, with slope $-2.069$ ($R^2=0.99898$).

\emph{Entanglement contour.}  The strip entanglement is resolved with the entanglement contour \cite{ChenVidal2014}.  A procedural convention matters and is used throughout this work: the subsystem is always the \emph{full-width} strip in $x$; any $x$-window is applied only to the already-computed contour sum, as a diagnostic restriction.  (Redefining the subsystem would change the entanglement problem.)  With this convention, the full-width integrated contour follows the log-chord law with slope $0.3439$, and three-column windows centered on each wall give $0.1686$ per wall: $0.3439\simeq2\times0.1686$, the cleanest statement that the strip entanglement is the sum of two equal chiral wall contributions, each with $c_{\rm wall}/3=1/6$.

\emph{Flattened-Hamiltonian surrogate.}  The OW flattened Hamiltonian $H_{\rm flat}=\sum_{\bm r,\nu}\big(P_{\nu+}(\bm r)-P_{\nu-}(\bm r)\big)$, with target covariance $\GOW=2P_{\rm occ}-\mathbf{1}$, reproduces the exact contour slopes closely (full width $0.3448$, wall windows $0.1691$ for $X/Y$ trial orbitals), validating the OW construction as the measurement target underlying the adaptive circuit.

\begin{figure}[h]
  \figplaceholder{1.6in}{%
  \textbf{FIG.~S3 PLACEHOLDER (exact benchmark, $1\times3$ panels).}\\[2pt]
  (a) Squared-correlator log-chord fits per column $x$ (slopes $-2.02$ at walls, steep interior decay); (b) integrated contour vs.\ log-chord coordinate, full-width and per-wall windows; (c) exact vs.\ flattened-Hamiltonian ($X/Y$ trial) contour slopes.\\
  {[}data: \texttt{exact\_DW\_benchmark\_notes}, \texttt{notebooks/exact\_DW\_GS\_benchmark.ipynb}{]}}
  \caption{Exact domain-wall ground-state benchmark.}
  \label{fig:sm-exact}
\end{figure}

%=====================================================================
\section{Overcomplete-Wannier truncation: form factors, hard vs.\ soft walls, and the soft-wall verification}
\label{sm:owtrunc}

\subsection{Finite-shell truncation and the effective form factor}

The topological obstruction relevant to the measurement circuit is carried not by the truncated OW spinor itself but by its \emph{overlap form factor} with the target band, $f_{\nu,n}(\bm k)\propto\langle\psi_n(\bm k)|\tau_\nu\rangle$, whose zero and winding encode the Chern number \cite{BhuiyanPanJian2026}.  A finite real-space mask $g(\bm r)$ dresses this overlap into the effective form factor
\begin{equation}
  f^{(g)}(\bm k')=\frac{1}{N}\sum_{\bm k\in{\rm BZ}}
  \tilde g(\bm k'-\bm k)\,e^{-\mathrm{i}(\bm k'-\bm k)\cdot\bm r}\,
  f(\bm k)\,\braket{u_{\bm k'}}{u_{\bm k}},
  \label{eq:sm-formfactor}
\end{equation}
a Bloch-overlap-dressed convolution.  The truncation criterion is then sharp: the finite-range circuit retains the Chern obstruction if and only if the zero and winding of $f^{(g)}$ survive the convolution.  For the production case ($\alpha=1$, square $\nsh$ mask), the OW spread scales are $\xi_{\rm rms}=0.674a$ and $\xi_{\rm exp}=1.443a$; the $\nsh=1$ mask retains $0.9922$ of the OW norm and preserves the zero, displaced from $k_x/\pi=0.5$ to $0.4922$ ($\nsh=2$: norm $0.9991$).  The $\nsh=0$ (on-site) limit instead performs a full Brillouin-zone average and destroys the controlled zero.  We caution that a naive Chern number assigned to the compact-support spinor is meaningless here: a smooth, nowhere-vanishing compact-support section has vanishing FHS Chern number even while the overlap winding remains $\nu=-1$; the form factor, not the spinor, carries the topology.

\subsection{Two domain-wall constructions: soft walls and hard walls}

The Letter's two equivalent adaptive constructions of the interface correspond to two OW-support conventions:
\begin{itemize}
\item \emph{Soft wall} (\texttt{dw\_truncation} off): measure OW occupancies of distinct target topologies on complementary regions ($\atopo$ slab, $\atriv$ exterior); OW modes near the wall retain exponentially small tails across it.
\item \emph{Hard wall} (\texttt{dw\_truncation} on): additionally mask each OW mode to the sector containing its center and renormalize.  This is the projector-support limit of an infinite-mass trivial exterior---an \emph{active boundary condition}, not a numerical convenience.
\end{itemize}
Two facts support the hard-wall interpretation.  First, the half-filled Chern marker of the flattened-Hamiltonian state: without sector truncation the marker saturates near $0.795$--$0.796$ (at $\nsh=1$ and even at full support), whereas full-support sector truncation restores $\calC=0.99999895$; entanglement is far more forgiving, with the $\nsh=1$ contour slopes already at their chiral values ($0.167$/$0.334$).  Second, the trivial-mass sweep at fixed topological interior: scanning $\atriv$ with hard walls, the extracted interface central charge counts the exposed chiral channels, $c_{\rm est}\approx2$ while the exterior is still topological, dropping to a robust $c_{\rm est}\approx1$ once the exterior crosses its transition near $\alpha=2$.  Hard-wall truncation exposes interface channels; it does not create chirality by itself.

\subsection{Soft-wall verification of the main-text results}
\label{sm:dwtrunc0}

The main text uses the hard-wall construction; here we verify that the soft-wall construction yields the same physics (Fig.~\ref{fig:sm-dwtrunc0}).  In the small-system maximally mixed purification campaign at $16\times32$, $100$ cycles, \emph{without} sector truncation, perfect correction purifies essentially completely---final total entropy $4.4\times10^{-4}$ ($\nsh=1$) and $1.3\times10^{-3}$ ($\nsh=2$) in natural-log units---while deterministic postselection saturates at $2\log2$ ($1.3862869$ and $1.3862943$, respectively).  The same postselected $2\log2$ residual recurs across truncation conventions and geometries, while the perfect-correction residual shows quantitative (not qualitative) sensitivity to truncation and geometry: $6.3\times10^{-4}$ bits (soft wall) vs.\ $2.8\times10^{-2}$ bits (hard wall) at $16\times32$, $\nsh=1$, and $0.234$ bits for the larger $30\times40$ hard-wall comparison at the same depth.  Entropy-slope windows and correlator diagnostics behave analogously across the two constructions.  We also note that the steered ensemble approaches but does not exactly reach the flattened-Hamiltonian covariance: in the postselected control at $16\times20$, the normalized Frobenius distance to $\GOW$ relaxes rapidly and saturates at $0.75$ (soft wall) and $0.72$ (hard wall) by $40$ cycles.  The adaptive steady state is a stochastic ensemble organized around the OW target, not an exact copy of it.

\begin{figure}[h]
  \figplaceholder{1.8in}{%
  \textbf{FIG.~S4 PLACEHOLDER (soft-wall verification, $2\times2$ panels).}\\[2pt]
  (a) Maximally mixed purification at $16\times32$, soft wall: total entropy vs.\ cycle, PC vs.\ PS ($2\log2$ line); (b) same with hard wall (comparison); (c) entropy-slope windows, soft vs.\ hard wall; (d) trivial-mass sweep: $c_{\rm est}$ vs.\ $\atriv$ with hard wall, showing the $2\to1$ chiral-channel step at the exterior transition.\\
  {[}data: \texttt{colab\_small\_system\_testing} purification summaries; \texttt{flattened\_hamiltonian\_analysis} $\alpha$-sweep{]}}
  \caption{Soft-wall (\texttt{dw\_truncation} off) verification and chiral-channel counting.}
  \label{fig:sm-dwtrunc0}
\end{figure}

%=====================================================================
\section{Correlator fit windows and entropy slope excess}
\label{sm:corr}

\subsection{Window-controlled correlator exponents}

The squared-correlator exponent $\Delta_{\rm fit}$ (minus the log-log slope vs.\ chord distance) is a crossover-sensitive diagnostic at accessible sizes.  In the streaming characterization ($20\times\{30,40,50\}$, $100$ PC samples, final cycle $2N_y$), default long-distance windows give unstable effective exponents---$1.41$, $1.57$, $1.63$ for perfect correction; $1.62$, $3.88$, $2.70$ for the single-sample postselected control---and even the \emph{equilibrium} flattened-Hamiltonian reference state gives default-window exponents of $4.9$--$6.4$ with poor fit quality ($R^2\simeq0.73$--$0.75$): the default window itself is contaminated by finite-size curvature.  On controlled intermediate windows $r\in[2,N_y/4]$, by contrast, perfect correction gives $2.130$, $2.117$, $2.119$ and the flattened reference gives $2.226$, $2.236$, $2.197$, both consistent with the chiral value $\Delta=2$.  The internal decision-table analysis quantifies the window dominance (maximal window-to-window spread $1.95$ at $N_x=20$, far above the $0.5$ stability threshold) and finds that postselection is \emph{not} uniformly cleaner (spread $2.26$ across $N_y$).  Correlator exponents should accordingly be quoted only with their windows, as consistency checks rather than precision universality data.

\subsection{Entropy slope excess}

The final-cycle entropy slopes of the perfect-correction ensemble exceed $1/3$ slightly: $m_s=0.3598$, $0.3492$, $0.3479$ for $N_y=30,40,50$ at cycles $60,80,100$ ($100$ trajectories), i.e., $\delta m_s=0.0265$, $0.0159$, $0.0146$.  Diagnostic tests on the saved observables resolve its origin:
(i) \emph{global charge imbalance is not the cause}---conditioning on low total-charge subsets does not remove the excess, and pooled correlations between charge and slope are weak;
(ii) \emph{local wall metrics are implicated}---conditioning on low wall-charge RMS, low wall frustration, or low Chern-marker error removes a substantial fraction of the excess, increasingly so at larger $N_y$;
(iii) \emph{incomplete convergence is implicated}---cycle-convergence fits extrapolate to smaller plateau excesses, most strongly for $N_y=40,50$;
(iv) \emph{averaging order is excluded}---the mean of per-sample slopes and the slope of the mean entropy curve agree to numerical roundoff.
The safe statement of the Letter follows: slopes near $1/3$ with finite-size and finite-depth corrections attributable to local wall defects and finite circuit depth.  In the time-scaled large-$N_y$ campaign ($T=N_y/2$, $10$ samples), the slope decreases monotonically toward $1/3$ through $N_y=100$ ($0.3615\to0.3470$); the largest point $N_y=120$ ($0.3569$) is non-monotone, consistent with its reduced effective depth and sample count, and we therefore quote the trend through $N_y=100$ in the main text and display $N_y=120$ here.

\begin{figure}[h]
  \figplaceholder{1.8in}{%
  \textbf{FIG.~S5 PLACEHOLDER (fit diagnostics, $2\times2$ panels).}\\[2pt]
  (a) Correlator exponent vs.\ fit window (PC, PS, flattened reference) showing window dominance; (b) selected-window exponents vs.\ $N_y$ with $\Delta=2$ line; (c) slope excess $\delta m_s$ vs.\ cycle with convergence fits, including the non-monotone $N_y=120$ point; (d) conditioned-subset slopes (low wall charge/frustration/Chern error) vs.\ unconditioned.\\
  {[}data: \texttt{correlator\_scaling\_diagnostics}, \texttt{perfect\_correction\_slope\_excess\_diagnostics} tables{]}}
  \caption{Correlator-window and slope-excess diagnostics.}
  \label{fig:sm-corr}
\end{figure}

%=====================================================================
\section{Modular charge spreading}
\label{sm:modular}

The chirality diagnostic of the Letter evolves localized charge packets under the modular Hamiltonian of the prepared state.  On the $20\times40$ torus with walls at $x=6,14$, the subsystem is a half cylinder ($20\times20$, $800$ single-particle modes).  To isolate the translation-invariant entanglement structure, the reduced covariance is averaged over all $y$-translated half-cuts before constructing the modular kernel
\begin{equation}
  h_{\rm mod}=-2\,V\,\mathrm{diag}\big(\mathrm{arctanh}\,g_i\big)\,V^\dagger,
  \qquad
  \bar G_A=V\,\mathrm{diag}(g_i)\,V^\dagger .
  \label{eq:sm-modular}
\end{equation}
Charge is injected on the wall columns (or their $\pm1$ neighborhoods) at the entanglement-cut rows and evolved by exact spectral modular evolution ($\Delta t=0.05$, $T=40$); total charge is conserved to better than $5\times10^{-7}$ in all runs.

Three findings support and sharpen the main-text panel.  First, packets launched on opposite walls acquire \emph{oppositely signed} center-of-mass drifts along the wall direction, as required for counter-propagating chiral interfaces; the averaged displacement signs are unambiguous even where single-launch velocity fits have low $R^2$, which is why the Letter quotes signed-drift evidence rather than a velocity.  Second, the $y$-averaged half-system \emph{modular spectra} of the exact domain-wall ground state and all eight flattened-OW surrogate variants coincide to all reported digits ($\varepsilon_{\rm max}=-\varepsilon_{\rm min}=23.7190$, $800$ modes, zero Hermiticity error): at the level of the translation-averaged entanglement spectrum, the OW targets are indistinguishable from the exact state.  Third, spatially resolved modular observables \emph{do} distinguish the constructions: by $T=40$ the topological-region charge fraction is $0.959$ for the exact state, $0.896$/$0.859$ for soft-wall OW surrogates ($\nsh=1,2$), and exactly $1.000$ for hard-wall (sector-truncated) variants, which impose exact confinement by construction.  Spectral agreement does not imply eigenvector agreement; this is the modular-language restatement of the soft/hard-wall distinction of Sec.~\ref{sm:owtrunc}, and the dynamical-transport counterpart of static modular-data probes of chiral central charge \cite{KimShiKatoAlbert2022,KimShiKatoAlbert2022b,LiHaldane2008}.

\begin{figure}[h]
  \figplaceholder{1.8in}{%
  \textbf{FIG.~S6 PLACEHOLDER (modular diagnostics, $2\times2$ panels).}\\[2pt]
  (a) Modular spectra of exact and OW-variant states (coinciding); (b) modular-time charge maps at $t=0,20,40$ showing wall-guided chiral fanning; (c) center-of-mass drift $\Delta y(t)$ per wall and per injection config; (d) topological-region charge fraction vs.\ time for exact / soft-wall / hard-wall states.\\
  {[}data: \texttt{notebooks/modular\_charge\_spreading} CSV summaries; \texttt{colab\_small\_system\_testing} dynamic modular spreading{]}}
  \caption{Modular charge spreading: spectra, chirality, and confinement.}
  \label{fig:sm-modular}
\end{figure}

%=====================================================================
\section{Transfer and stability diagnostics}
\label{sm:stability}

\subsection{Linearized update and raster contraction}

Linearizing a single rank-one OW measurement-feedback event about its own fixed branch ($G_\ast P=PG_\ast=sP$, $s=\pm1$, $P=\ket{W}\bra{W}$) gives
\begin{equation}
  \delta G'=(\mathbf{1}-P)\,\delta G\,(\mathbf{1}-P),
  \label{eq:sm-rank1}
\end{equation}
independently of the branch sign; a full raster cycle composes these complementary projectors as $\delta G\mapsto L\,\delta G\,L^\dagger$ with $L=Q_m\cdots Q_1$, $Q_i=\mathbf{1}-P_i$.  All inspected cases (including finite-shell, hard-wall geometries) have $\|L\|<1$, the analytic mechanism behind the $\mathcal{O}(1)$ contraction of the bulk.  This local statement does not by itself fix the wall sector, which is probed numerically below.

\subsection{Finite-time Lyapunov spectra of the single-particle transfer frame}

Along realized trajectories we propagate a frame $V\in\mathbb{C}^{2N_xN_y\times n_{\rm vec}}$ in the top-layer single-particle space ($n_{\rm vec}=2N_xN_y$), QR-reorthogonalized once per cycle, and accumulate per-cycle-normalized finite-time Lyapunov exponents $\lambda_i(T)$, with gap proxy $\Dlam(T)=\min_i|\lambda_i(T)|$.  We emphasize the diagnostic's scope: this is the \emph{single-particle transfer-frame} spectrum of the stochastic dynamics, not the full tangent spectrum on Hermitian covariance perturbations; the two are related but distinct linearizations.  The campaign uses $N_x=20$, $N_y=30,40,50$, $T=N_y$, $25$ samples (PC) and $1$ (PS).  Table~\ref{tab:lyap} collects the final-time gaps and Chern markers.

\begin{table}[h]
\caption{\label{tab:lyap}Final-time single-particle Lyapunov gaps $\Dlam$ and real-space Chern markers $\calC$.  Uniform geometry (no wall) vs.\ domain-wall geometry; perfect correction (PC, $25$ samples, mean(spread)) and postselection (PS, single deterministic trajectory).}
\begin{ruledtabular}
\begin{tabular}{llccc}
 & & $N_y=30$ & $N_y=40$ & $N_y=50$ \\
\colrule
uniform & $\Dlam$ (PC) & $0.722(11)$ & $0.748(8)$ & $0.768(9)$ \\
        & $\Dlam$ (PS) & $0.945$ & $0.955$ & $0.952$ \\
        & $\calC$       & \multicolumn{3}{c}{quantized, $\simeq1$} \\
\colrule
domain wall & $\Dlam$ (PC) & $0.027$ & $0.024$ & $0.020$ \\
            & $\Dlam$ (PS) & $0.021$ & $0.0076$ & $0.0078$ \\
            & $\calC$ (PC)  & $0.978(6)$ & $0.954(16)$ & $0.954(17)$ \\
            & $\calC$ (PS)  & $0.662$ & $0.863$ & $0.883$ \\
\end{tabular}
\end{ruledtabular}
\end{table}

The uniform gap plateaus within ${\sim}10$ cycles and is size-independent; the domain-wall gap is suppressed by one to two orders of magnitude, decreases monotonically with $N_y$ under perfect correction, and has \emph{not} plateaued by $T=N_y$.  The minimum-gap tangent vector is sharply weighted on the wall columns and extended along the wall, and per-sample correlations link smaller $\Dlam$ to smaller Chern markers.  Three circumferences at finite time cannot distinguish $1/N_y$, another power, or logarithmic closing; the thermodynamic scaling form is left open.

\subsection{Regularized Choi particle-transfer gap}

An independent transfer diagnostic tracks the cumulative two-leg Choi covariance $\Sigma$ of the realized circuit operator (legs $L,R$ label the Choi doubling, not spatial regions).  Each elementary branch contributes a formally singular rank-one tensor
\begin{equation}
  K(\eta_1,\eta_2)=
  \begin{pmatrix}
    \eta_1 P & Q\\
    Q & \eta_2 P
  \end{pmatrix},
  \qquad P=\ket{\chi}\bra{\chi},\quad Q=\mathbf{1}-P,
\end{equation}
whose contraction is regulated by $K^{(\epsilon)}_{LL}=\eta_1P+\epsilon Q$ and the Schur-complement limit $\epsilon\to0$; the update is finite whenever the denominator $d=\bra{\chi}\Sigma_{RR}\ket{\chi}-\eta_1$ is nonzero, and is implemented as a resolvent-based rank-one update with monitored $|d|$.  The finite-sector particle-transfer operator and gap are
\begin{equation}
  T_{\rm p}=-(\mathbf{1}+\Sigma_{LL})^{-1}\Sigma_{LR},
  \qquad
  \Dp(t)=\min_j\Big|\frac{\log s_j(T_{\rm p})}{t}\Big|,
  \label{eq:sm-choi}
\end{equation}
with $\Sigma_{LL}$ eigenvalues $a\in(-1,1)$ mapped to exponents $\lambda_{\rm p}(a;t)=[\log(1-a)-\log(1+a)]/2t$ and endpoint eigenvalues $a=\pm1$ treated as exact particle zeros/poles.  Two pipelines are used.  The GPU pipeline gives an early-cycle chart: domain-wall slab gaps sit at $\Dp\sim(1$--$2)\times10^{-2}$ while the uniform control shows much larger transient gaps; all samples are censored once the Choi chart leaves its accepted numerical domain (stable fraction reaching zero by cycles $15$--$16$), so no late-time GPU gap is quoted.  The CPU pipeline reaches final cycles with exact $\Sigma_{LL}$ eigensolves: the production campaign ($N_x=20$, $N_y=20,30,40$, $\atopo\in\{1,3\}$, $\atriv=30$, $10$ samples, $2N_y$ cycles, slab-restricted, rank-one updates with no dense fallback) completes with ten stable final samples per configuration and saved final spectra and eigenvectors.  Endpoint saturation is the caveat to track: cycles may carry fewer finite transfer exponents than requested (NaN-padded), and the $\atopo=3$ runs saturate to endpoint sentinels by cycles $20$--$40$---evidence of finite-sector collapse, not of a resolved finite gap.  The Lyapunov and Choi spectra are distinct operators; their qualitative agreement (a near-marginal wall channel atop a strongly contracting bulk), together with the wall-localized eigenvectors of both, is what the Letter reports.

\begin{figure}[h]
  \figplaceholder{1.8in}{%
  \textbf{FIG.~S7 PLACEHOLDER (transfer diagnostics, $2\times2$ panels).}\\[2pt]
  (a) $\Dlam(t)$ vs.\ cycle: uniform plateau vs.\ domain-wall non-plateau; (b) Choi stable fraction and $\Dp$ vs.\ cycle with censoring window (GPU) and final-cycle CPU points; (c) CPU final Choi spectra (rooted singular values) for $\atopo=1$ vs.\ endpoint-saturated $\atopo=3$; (d) near-gap Choi eigenvector density maps with wall locators.\\
  {[}data: \texttt{colab\_lyapunov}, \texttt{colab\_regularized\_choi\_transfer\_matrix}, \texttt{choi\_covariance\_cpu}{]}}
  \caption{Transfer-stability diagnostics: Lyapunov frame and Choi particle transfer.}
  \label{fig:sm-choi}
\end{figure}

%=====================================================================
\section{Circuit-depth bounds and the dimensional loophole}
\label{sm:depth}

\subsection{The entanglement-growth bound of local adaptive circuits}

Lu, Lessa, Kim, and Hsieh \cite{LuLessaKimHsieh2022} prove the following (their Appendix~D).  Consider a depth-$D$ adaptive circuit in which each layer consists of non-overlapping local gates, each a unitary or a measurement projector, with arbitrary outcome-dependent adaptivity.  Let $S_0=\log\chi$ be the max-entropy (Schmidt rank $\chi$) of a region $A$.  Then, \emph{independently of the measurement outcomes},
\begin{equation}
  S_0^{\rm out}(A)\;\leq\;S_0^{\rm in}(A)+\mathcal{O}(D\,|\partial A|),
  \label{eq:sm-bound}
\end{equation}
because (i) a gate supported entirely in $A$ or in $\bar A$ cannot increase the Schmidt rank, and (ii) a gate straddling the cut, having $\mathcal{O}(1)$ support, multiplies the Schmidt rank by at most an $\mathcal{O}(1)$ factor; each layer contains $\mathcal{O}(|\partial A|)$ straddling gates.  Since $S_n\leq S_0$ for all R\'enyi indices $n>0$, a circuit acting on a \emph{product state} outputs states with $S_n(A)\leq\mathcal{O}(D|\partial A|)$ for every region.  (The product-state input matters: for general inputs the von Neumann entropy alone admits no such bound under measurement.)  In $d=1$, intervals have $|\partial A|=\mathcal{O}(1)$, while a CFT ground state of size $L$ has $S\sim\tfrac{c}{3}\log L$; hence \emph{any} local adaptive preparation of a 1d critical state requires
\begin{equation}
  D\;\gtrsim\;\mathcal{O}(\log L),
  \label{eq:sm-loglower}
\end{equation}
and the adaptive-MERA construction of Ref.~\cite{LuLessaKimHsieh2022} saturates it.

\subsection{Why the bound is not binding for interface criticality}

Our circuit satisfies every hypothesis of Eq.~\eqref{eq:sm-bound}---local gates, local measurements, product (or maximally mixed) inputs, no classical communication at all---but acts on a two-dimensional lattice, and the critical degrees of freedom live on a one-dimensional \emph{interface} of the dynamics.  For any region probing a wall interval of length $\ell$ (a strip or rectangle of width $w$), the boundary is $|\partial A|=\mathcal{O}(\ell+w)$, so a depth-$\mathcal{O}(1)$ circuit has entanglement budget $\mathcal{O}(\ell+w)\gg\tfrac13\log\ell$.  Equation~\eqref{eq:sm-bound} therefore permits constant-depth interface criticality for \emph{every} region simultaneously; the $\Omega(\log L)$ conclusion (\ref{eq:sm-loglower}) follows from the bound only when the circuit is local in the \emph{same} dimension as the critical state, i.e., when cuts have $\mathcal{O}(1)$ boundary.  That is the single assumption our construction sidesteps.  The dynamical bulk--boundary correspondence \cite{BhuiyanPanJian2026} then shows the permitted scaling is actually \emph{achieved}: the effective-Lindbladian convergence time of the bulk protocol is $T_{\rm conv}\sim\mathcal{O}(1)$, and the trajectory data of the Letter show chiral log-law entanglement already at depths of order tens of cycles, each of $\mathcal{O}(1)$ circuit depth at fixed $\nsh$.

\subsection{Honest accounting and relation to other measurement shortcuts}

Three costs and qualifications accompany this statement.  (i) \emph{Spatial overhead}: an $L\times L$ bulk hosts an $\mathcal{O}(L)$ interface; the depth saving is bought with quadratically many auxiliary degrees of freedom.  This is the dynamical counterpart of preparing critical \emph{boundary} states by measuring two-dimensional resource states \cite{TTVV2024,Zhu2023,LeeJiBiFisher2022,ZhuBoundary2023}, with two differences: no resource state is consumed (the same measurements continuously re-stabilize the ensemble), and no outcome decoding is required (feedback corrects rather than interprets the record).  (ii) \emph{Ensemble preparation}: the protocol prepares a stochastic ensemble, every member of which is pure and critical at the interface; it does not output one predetermined wavefunction.  This differs from constant-depth preparations of critical \emph{mixed} states, which use nonlocal classical communication and realize criticality in negativity atop volume-law entropy \cite{LuZhangVijayHsieh2023,LuMixedFollowup2025}; here criticality is carried by the entanglement entropy of each pure trajectory.  (iii) \emph{Interface-only criticality}: the bulk remains gapped and area-law---which is precisely why the no-go premise is inapplicable rather than violated.  Finally, since the protocol uses no classical communication, it cannot exploit measurement-enabled information speedups \cite{Friedman2022}; the logarithmic entanglement is generated by genuinely local dynamics in one higher dimension.

%=====================================================================
\begin{thebibliography}{99}

\bibitem{BhuiyanPanJian2026}
A.~Bhuiyan, H.~Pan, and C.-M. Jian,
``Free-fermion dynamics with measurements: Topological classification and adaptive preparation of topological states,''
Phys. Rev. Res. \textbf{8}, 023147 (2026), arXiv:2507.13437.

\bibitem{Rashba1997}
E.~I. Rashba, L.~E. Zhukov, and A.~L. Efros,
``Orthogonal localized wave functions of an electron in a magnetic field,''
Phys. Rev. B \textbf{55}, 5306 (1997).

\bibitem{Qi2011}
X.-L. Qi,
``Generic wave-function description of fractional quantum anomalous Hall states and fractional topological insulators,''
Phys. Rev. Lett. \textbf{107}, 126803 (2011).

\bibitem{JianQi2013}
C.-M. Jian and X.-L. Qi,
``Crystal-symmetry preserving Wannier states for fractional Chern insulators,''
Phys. Rev. B \textbf{88}, 165134 (2013).

\bibitem{LiZou2023}
Y.~Li, Y.~Zou, P.~Glorioso, E.~Altman, and M.~P.~A. Fisher,
``Cross entropy benchmark for measurement-induced phase transitions,''
Phys. Rev. Lett. \textbf{130}, 220404 (2023).

\bibitem{GarrattAltman2024}
S.~J. Garratt and E.~Altman,
``Probing postmeasurement entanglement without postselection,''
PRX Quantum \textbf{5}, 030311 (2024).

\bibitem{Guo2025}
J.~Guo, O.~Hart, C.-F. Chen, A.~J. Friedman, and A.~Lucas,
``Designing open quantum systems with known steady states: Davies generators and beyond,''
Quantum \textbf{9}, 1612 (2025).

\bibitem{Dolgirev2020}
P.~E. Dolgirev, J.~Marino, D.~Sels, and E.~Demler,
``Non-Gaussian correlations imprinted by local dephasing in fermionic wires,''
Phys. Rev. B \textbf{102}, 100301(R) (2020).

\bibitem{BarthelZhang2022}
T.~Barthel and Y.~Zhang,
``Solving quasi-free and quadratic Lindblad master equations for open fermionic and bosonic systems,''
J. Stat. Mech. (2022) 113101.

\bibitem{LieuMcGinleyCooper2020}
S.~Lieu, M.~McGinley, and N.~R. Cooper,
``Tenfold way for quadratic Lindbladians,''
Phys. Rev. Lett. \textbf{124}, 040401 (2020).

\bibitem{AltlandFleischhauerDiehl2021}
A.~Altland, M.~Fleischhauer, and S.~Diehl,
``Symmetry classes of open fermionic quantum matter,''
Phys. Rev. X \textbf{11}, 021037 (2021).

\bibitem{Goldstein2019}
M.~Goldstein,
``Dissipation-induced topological insulators: A no-go theorem and a recipe,''
SciPost Phys. \textbf{7}, 067 (2019).

\bibitem{BudichDiehl2015}
J.~C. Budich and S.~Diehl,
``Topology of density matrices,''
Phys. Rev. B \textbf{91}, 165140 (2015).

\bibitem{Bardyn2018}
C.-E. Bardyn, L.~Wawer, A.~Altland, M.~Fleischhauer, and S.~Diehl,
``Probing the topology of density matrices,''
Phys. Rev. X \textbf{8}, 011035 (2018).

\bibitem{WawerFleischhauer2021}
L.~Wawer and M.~Fleischhauer,
``Chern number and Berry curvature for Gaussian mixed states of fermions,''
Phys. Rev. B \textbf{104}, 094104 (2021).

\bibitem{Diehl2011}
S.~Diehl, E.~Rico, M.~A. Baranov, and P.~Zoller,
``Topology by dissipation in atomic quantum wires,''
Nat. Phys. \textbf{7}, 971 (2011).

\bibitem{Bardyn2013}
C.-E. Bardyn, M.~A. Baranov, C.~V. Kraus, E.~Rico, A.~\.{I}mamo\u{g}lu, P.~Zoller, and S.~Diehl,
``Topology by dissipation,''
New J. Phys. \textbf{15}, 085001 (2013).

\bibitem{BudichZollerDiehl2015}
J.~C. Budich, P.~Zoller, and S.~Diehl,
``Dissipative preparation of Chern insulators,''
Phys. Rev. A \textbf{91}, 042117 (2015).

\bibitem{YangMoligniniBergholtz2023}
F.~Yang, P.~Molignini, and E.~J. Bergholtz,
``Dissipative boundary state preparation,''
Phys. Rev. Res. \textbf{5}, 043229 (2023).

\bibitem{HegdeEhmckeMeng2023}
S.~S. Hegde, T.~Ehmcke, and T.~Meng,
``Edge-selective extremal damping from topological heritage of dissipative Chern insulators,''
Phys. Rev. Lett. \textbf{131}, 256601 (2023).

\bibitem{ShavitGoldstein2020}
G.~Shavit and M.~Goldstein,
``Topology by dissipation: Transport properties,''
Phys. Rev. B \textbf{101}, 125412 (2020).

\bibitem{ChenVidal2014}
Y.~Chen and G.~Vidal,
``Entanglement contour,''
J. Stat. Mech. (2014) P10011.

\bibitem{KimShiKatoAlbert2022}
I.~H. Kim, B.~Shi, K.~Kato, and V.~V. Albert,
``Chiral central charge from a single bulk wave function,''
Phys. Rev. Lett. \textbf{128}, 176402 (2022).

\bibitem{KimShiKatoAlbert2022b}
I.~H. Kim, B.~Shi, K.~Kato, and V.~V. Albert,
``Modular commutator in gapped quantum many-body systems,''
Phys. Rev. B \textbf{106}, 075147 (2022).

\bibitem{LiHaldane2008}
H.~Li and F.~D.~M. Haldane,
``Entanglement spectrum as a generalization of entanglement entropy,''
Phys. Rev. Lett. \textbf{101}, 010504 (2008).

\bibitem{LuLessaKimHsieh2022}
T.-C. Lu, L.~A. Lessa, I.~H. Kim, and T.~H. Hsieh,
``Measurement as a shortcut to long-range entangled quantum matter,''
PRX Quantum \textbf{3}, 040337 (2022).

\bibitem{TTVV2024}
N.~Tantivasadakarn, R.~Thorngren, A.~Vishwanath, and R.~Verresen,
``Long-range entanglement from measuring symmetry-protected topological phases,''
Phys. Rev. X \textbf{14}, 021040 (2024).

\bibitem{Zhu2023}
G.-Y. Zhu, N.~Tantivasadakarn, A.~Vishwanath, S.~Trebst, and R.~Verresen,
``Nishimori's cat: Stable long-range entanglement from finite-depth unitaries and weak measurements,''
Phys. Rev. Lett. \textbf{131}, 200201 (2023).

\bibitem{LeeJiBiFisher2022}
J.~Y. Lee, W.~Ji, Z.~Bi, and M.~P.~A. Fisher,
``Decoding measurement-prepared quantum phases and transitions: from Ising model to gauge theory, and beyond,''
arXiv:2208.11699.

\bibitem{ZhuBoundary2023}
Y.~Wu, G.-Y. Zhu, \emph{et al.},
``Triggering boundary phase transitions through bulk measurements in 2D cluster states,''
arXiv:2305.14231.

\bibitem{LuZhangVijayHsieh2023}
T.-C. Lu, Z.~Zhang, S.~Vijay, and T.~H. Hsieh,
``Mixed-state long-range order and criticality from measurement and feedback,''
PRX Quantum \textbf{4}, 030318 (2023).

\bibitem{LuMixedFollowup2025}
Z.~Zhang, T.-C. Lu, S.~Vijay, and T.~H. Hsieh,
``Universal properties of critical mixed states from measurement and feedback,''
arXiv:2503.09597.

\bibitem{Friedman2022}
A.~J. Friedman, C.~Yin, Y.~Hong, and A.~Lucas,
``Locality and error correction in quantum dynamics with measurement,''
arXiv:2206.09929.

\end{thebibliography}

\end{document}
